%% kapa: exam=1 round=33 year=2022 period=1 form=A subject=civil_law subject_ko=민법 seq=2 number=2
%% source: Q-Net 감정평가사 기출문제 (cst00309) artlSeq=5209997 file=raw/1008/5209997/2022년도 제33회 감정평가사 1차시험 1교시(민법,경제학,부동산학).pdf
%% conversion: kordoc 4.24.51 -> tools/gen1.py (auto); review=proofread-sonnet; flags=none
\begin{question}{2}
신의칙에 관한 설명으로 옳은 것을 모두 고른 것은? (다툼이 있으면 판례에 따름)
\begin{framedtext}
ㄱ. 법원은 당사자의 주장이 없으면 직권으로 신의칙 위반 여부를 판단할 수 없다. \par
ㄴ. 무권대리인이 무권대리행위 후 단독으로 본인의 지위를 상속한 경우, 본인의 지위에서 그 무권대리행위의 추인을 거절하는 것은 신의칙에 반한다. \par
ㄷ. 부동산거래에서 신의칙상 고지의무의 대상은 직접적인 법령의 규정뿐만 아니라 계약상, 관습상 또는 조리상의 일반원칙에 의해서도 인정될 수 있다.
\end{framedtext}
\begin{choices}
  \item ㄱ
  \item ㄴ
  \item ㄱ, ㄷ
  \item ㄴ, ㄷ
  \item ㄱ, ㄴ, ㄷ
\end{choices}
\end{question}
\answer{4}
